% GENERATED FILE. Edit canonical lessons, then run npm run generate:books.
% canonical-source-hash: fnv1a64:c15b22e94517a9a1
% canonical-lessons: KA-C397-nagarika, KA-C397-pauratva, KA-C397-sarvajanika, KA-C397-samudaya, KA-C397-oggudu, KA-R397-review-391-397

\chapter{Citizens Together}
\label{ch:ka-citizens-together}
\hlchaptermodality{397}

\begin{chapteropening}
\textbf{By the end of this chapter:} \emph{I can say a citizen, citizenship, public, a community and to unite.}

Complete the last lesson of chapter 397: I can say a citizen, citizenship, public, a community and to unite.
\end{chapteropening}

\section[\texorpdfstring{\kn{ನಾಗರಿಕ} (nāgarika)}{ನಾಗರಿಕ (nāgarika)}]{\kn{ನಾಗರಿಕ} (nāgarika) — a citizen}

\label{lesson:KA-C397-nagarika}



\begin{quote}
Before the new one: say the Kannada for a legislator, then the Kannada for to endorse.
\end{quote}

\subsection*{You'll want to know: \kn{ನಾಗರಿಕ}}
\textbf{\kn{ನಾಗರಿಕ}} — \emph{nāgarika} — \textquotedblleft{}a citizen\textquotedblright{}.

\begin{quote}
\kn{ಪ್ರತಿಯೊಬ್ಬ} \kn{ನಾಗರಿಕನೂ} \kn{ತೆರಿಗೆ} \kn{ಕಟ್ಟಬೇಕು}.
\end{quote}

\begin{quote}
\emph{pratiyobba nāgarikanū terige kaṭṭabēku.}
\end{quote}

\begin{quote}
Every citizen must pay taxes.
\end{quote}

Plural: \kn{ನಾಗರಿಕರು}, \emph{nāgarikaru}.

As an adjective it means civil, civilised.

\begin{cousinweb}
\textbf{\kn{ನಾಗರಿಕ}} (\emph{nāgarika}) is borrowed from Sanskrit \emph{nāgarika}, \textquotedblleft{}born or living in a town, civic\textquotedblright{} (Monier-Williams), from \emph{nagara}, \textquotedblleft{}a town\textquotedblright{}.
\end{cousinweb}

\begin{grammarlens}[title={What you've built}]
One more word for giving reasons and opinions.
\end{grammarlens}

\subsection*{Your turn}
\begin{itemize}
\raggedright
  \item \emph{Say it:} \emph{nāgarika}
  \item \emph{Say it:} \emph{nāgarika}, once more
  \item \emph{Say it:} \emph{pratiyobba nāgarikanū terige kaṭṭabēku.} — Every citizen must pay taxes.
  \item \emph{Recall:} say the Kannada for a legislator, then the Kannada for to endorse, then say \emph{nāgarika} again
\end{itemize}

\begin{tcolorbox}[breakable,colback=teal!4,colframe=teal!35!black,title={Before you move on}]
What does \kn{ನಾಗರಿಕ} mean? (\textquotedblleft{}A citizen\textquotedblright{}.) Say it once more.
\end{tcolorbox}
\section[\texorpdfstring{\kn{ಪೌರತ್ವ} (pauratva)}{ಪೌರತ್ವ (pauratva)}]{\kn{ಪೌರತ್ವ} (pauratva) — citizenship}

\label{lesson:KA-C397-pauratva}



\begin{quote}
Before the new one: say the Kannada for to endorse, then the Kannada for a citizen.
\end{quote}

\subsection*{You'll want to know: \kn{ಪೌರತ್ವ}}
\textbf{\kn{ಪೌರತ್ವ}} — \emph{pauratva} — \textquotedblleft{}citizenship\textquotedblright{}.

\begin{quote}
\kn{ಅವಳಿಗೆ} \kn{ಈ} \kn{ದೇಶದ} \kn{ಪೌರತ್ವ} \kn{ಸಿಕ್ಕಿತು}.
\end{quote}

\begin{quote}
\emph{avaḷige ī dēśada pauratva sikkitu.}
\end{quote}

\begin{quote}
She got the citizenship of this country.
\end{quote}

\kn{ಪೌರ} (citizen) + -\kn{ತ್ವ} (-ship).

\begin{cousinweb}
\textbf{\kn{ಪೌರತ್ವ}} (\emph{pauratva}) is built on Sanskrit \emph{paura}, \textquotedblleft{}a citizen\textquotedblright{} (CDIAL 8419), with -\emph{tva}, \textquotedblleft{}-ship\textquotedblright{}.
\end{cousinweb}

\begin{grammarlens}[title={What you've built}]
One more word for giving reasons and opinions.
\end{grammarlens}

\subsection*{Your turn}
\begin{itemize}
\raggedright
  \item \emph{Say it:} \emph{pauratva}
  \item \emph{Say it:} \emph{pauratva}, once more
  \item \emph{Say it:} \emph{avaḷige ī dēśada pauratva sikkitu.} — She got the citizenship of this country.
  \item \emph{Recall:} say the Kannada for to endorse, then the Kannada for a citizen, then say \emph{pauratva} again
\end{itemize}

\begin{tcolorbox}[breakable,colback=teal!4,colframe=teal!35!black,title={Before you move on}]
What does \kn{ಪೌರತ್ವ} mean? (\textquotedblleft{}Citizenship\textquotedblright{}.) Say it once more.
\end{tcolorbox}
\section[\texorpdfstring{\kn{ಸಾರ್ವಜನಿಕ} (sārvajanika)}{ಸಾರ್ವಜನಿಕ (sārvajanika)}]{\kn{ಸಾರ್ವಜನಿಕ} (sārvajanika) — public}

\label{lesson:KA-C397-sarvajanika}



\begin{quote}
Before the new one: say the Kannada for a citizen, then the Kannada for citizenship.
\end{quote}

\subsection*{You'll want to know: \kn{ಸಾರ್ವಜನಿಕ}}
\textbf{\kn{ಸಾರ್ವಜನಿಕ}} — \emph{sārvajanika} — \textquotedblleft{}public\textquotedblright{}.

\begin{quote}
\kn{ಇದು} \kn{ಸಾರ್ವಜನಿಕ} \kn{ಉದ್ಯಾನ}.
\end{quote}

\begin{quote}
\emph{idu sārvajanika udyāna.}
\end{quote}

\begin{quote}
This is a public park.
\end{quote}

\kn{ಸಾರ್ವಜನಿಕರು}, the public.

\begin{cousinweb}
\textbf{\kn{ಸಾರ್ವಜನಿಕ}} (\emph{sārvajanika}) is borrowed from Sanskrit \emph{sārvajanika}, \textquotedblleft{}belonging or suited to all people\textquotedblright{} (Monier-Williams).
\end{cousinweb}

\begin{grammarlens}[title={What you've built}]
One more word for giving reasons and opinions.
\end{grammarlens}

\subsection*{Your turn}
\begin{itemize}
\raggedright
  \item \emph{Say it:} \emph{sārvajanika}
  \item \emph{Say it:} \emph{sārvajanika}, once more
  \item \emph{Say it:} \emph{idu sārvajanika udyāna.} — This is a public park.
  \item \emph{Recall:} say the Kannada for a citizen, then the Kannada for citizenship, then say \emph{sārvajanika} again
\end{itemize}

\begin{tcolorbox}[breakable,colback=teal!4,colframe=teal!35!black,title={Before you move on}]
What does \kn{ಸಾರ್ವಜನಿಕ} mean? (\textquotedblleft{}Public\textquotedblright{}.) Say it once more.
\end{tcolorbox}
\section[\texorpdfstring{\kn{ಸಮುದಾಯ} (samudāya)}{ಸಮುದಾಯ (samudāya)}]{\kn{ಸಮುದಾಯ} (samudāya) — a community}

\label{lesson:KA-C397-samudaya}



\begin{quote}
Before the new one: say the Kannada for citizenship, then the Kannada for public.
\end{quote}

\subsection*{You'll want to know: \kn{ಸಮುದಾಯ}}
\textbf{\kn{ಸಮುದಾಯ}} — \emph{samudāya} — \textquotedblleft{}a community\textquotedblright{}.

\begin{quote}
\kn{ನಮ್ಮ} \kn{ಸಮುದಾಯ} \kn{ಒಟ್ಟಿಗೆ} \kn{ಹಬ್ಬ} \kn{ಆಚರಿಸುತ್ತದೆ}.
\end{quote}

\begin{quote}
\emph{namma samudāya oṭṭige habba ācarisuttade.}
\end{quote}

\begin{quote}
Our community celebrates the festival together.
\end{quote}

\kn{ಸಮುದಾಯ} \kn{ಭವನ}, a community hall.

\begin{cousinweb}
\textbf{\kn{ಸಮುದಾಯ}} (\emph{samudāya}) is borrowed from Sanskrit \emph{samudāya}, \textquotedblleft{}a combination, a multitude, a whole\textquotedblright{} (Monier-Williams).
\end{cousinweb}

\begin{grammarlens}[title={What you've built}]
One more word for giving reasons and opinions.
\end{grammarlens}

\subsection*{Your turn}
\begin{itemize}
\raggedright
  \item \emph{Say it:} \emph{samudāya}
  \item \emph{Say it:} \emph{samudāya}, once more
  \item \emph{Say it:} \emph{namma samudāya oṭṭige habba ācarisuttade.} — Our community celebrates the festival together.
  \item \emph{Recall:} say the Kannada for citizenship, then the Kannada for public, then say \emph{samudāya} again
\end{itemize}

\begin{tcolorbox}[breakable,colback=teal!4,colframe=teal!35!black,title={Before you move on}]
What does \kn{ಸಮುದಾಯ} mean? (\textquotedblleft{}A community\textquotedblright{}.) Say it once more.
\end{tcolorbox}
\section[\texorpdfstring{\kn{ಒಗ್ಗೂಡು} (oggūḍu)}{ಒಗ್ಗೂಡು (oggūḍu)}]{\kn{ಒಗ್ಗೂಡು} (oggūḍu) — to unite}

\label{lesson:KA-C397-oggudu}



\begin{quote}
Before the new one: say the Kannada for public, then the Kannada for a community.
\end{quote}

\subsection*{You'll want to know: \kn{ಒಗ್ಗೂಡು}}
\textbf{\kn{ಒಗ್ಗೂಡು}} — \emph{oggūḍu} — \textquotedblleft{}to unite\textquotedblright{}.

\begin{quote}
\kn{ಊರಿನ} \kn{ಜನರು} \kn{ಕೆರೆ} \kn{ಉಳಿಸಲು} \kn{ಒಗ್ಗೂಡಿದರು}.
\end{quote}

\begin{quote}
\emph{ūrina janaru kere uḷisalu oggūḍidaru.}
\end{quote}

\begin{quote}
The townspeople came together to save the lake.
\end{quote}

Past (I …): \kn{ಒಗ್ಗೂಡಿದೆ}, \emph{oggūḍide}.

\kn{ಒಗ್ಗೂಡಿಸು}, to unite others.

\begin{cousinweb}
\textbf{\kn{ಒಗ್ಗೂಡು}} (\emph{oggūḍu}) is native Dravidian, built from \emph{oggu}, \textquotedblleft{}to be one, to be united\textquotedblright{} (DEDR 990), and \emph{kūḍu}, \textquotedblleft{}to join, to come together\textquotedblright{} (DEDR 1882).
\end{cousinweb}

\begin{grammarlens}[title={What you've built}]
One more word for giving reasons and opinions.
\end{grammarlens}

\subsection*{Your turn}
\begin{itemize}
\raggedright
  \item \emph{Say it:} \emph{oggūḍu}
  \item \emph{Say it:} \emph{oggūḍu}, once more
  \item \emph{Say it:} \emph{ūrina janaru kere uḷisalu oggūḍidaru.} — The townspeople came together to save the lake.
  \item \emph{Recall:} say the Kannada for public, then the Kannada for a community, then say \emph{oggūḍu} again
\end{itemize}

\begin{tcolorbox}[breakable,colback=teal!4,colframe=teal!35!black,title={Before you move on}]
What does \kn{ಒಗ್ಗೂಡು} mean? (\textquotedblleft{}To unite\textquotedblright{}.) Say it once more.
\end{tcolorbox}
\section[(dialogue)]{Looking back: from Fairness for All to Citizens Together}

\label{lesson:KA-R397-review-391-397}



\begin{quote}
Say the Kannada for equality and to discriminate.
\end{quote}

\subsection*{Your turn}
Say these aloud:

\begin{itemize}
\raggedright
  \item equality, to discriminate, to exploit, empowerment, commitment
  \item an election, a vote, a majority, to campaign, to address a crowd
  \item a parliament, a minister, a constitution, democracy, to represent
  \item politics, a politician, a political party, to boycott, leadership
  \item administration, a ministry, a village council, to implement, a committee
  \item a draft law, to approve, an amendment, a legislator, to endorse
  \item a citizen, citizenship, public, a community, to unite
\end{itemize}

\begin{tcolorbox}[breakable,colback=teal!4,colframe=teal!35!black,title={Before you move on}]
Equality? (\textbf{\kn{ಸಮಾನತೆ}}, \emph{samānate}.) To discriminate? (\textbf{\kn{ತಾರತಮ್ಯ} \kn{ಮಾಡು}}, \emph{tāratamya māḍu}.) To exploit? (\textbf{\kn{ಶೋಷಿಸು}}, \emph{śōṣisu}.) Empowerment? (\textbf{\kn{ಸಬಲೀಕರಣ}}, \emph{sabalīkaraṇa}.) Commitment? (\textbf{\kn{ಬದ್ಧತೆ}}, \emph{baddhate}.) An election? (\textbf{\kn{ಚುನಾವಣೆ}}, \emph{cunāvaṇe}.) A vote? (\textbf{\kn{ಮತ}}, \emph{mata}.) A majority? (\textbf{\kn{ಬಹುಮತ}}, \emph{bahumata}.) To campaign? (\textbf{\kn{ಪ್ರಚಾರ} \kn{ಮಾಡು}}, \emph{pracāra māḍu}.) To address a crowd? (\textbf{\kn{ಸಂಬೋಧಿಸು}}, \emph{sambōdhisu}.) A parliament? (\textbf{\kn{ಸಂಸತ್ತು}}, \emph{saṃsattu}.) A minister? (\textbf{\kn{ಮಂತ್ರಿ}}, \emph{mantri}.) A constitution? (\textbf{\kn{ಸಂವಿಧಾನ}}, \emph{saṃvidhāna}.) Democracy? (\textbf{\kn{ಪ್ರಜಾಪ್ರಭುತ್ವ}}, \emph{prajāprabhutva}.) To represent? (\textbf{\kn{ಪ್ರತಿನಿಧಿಸು}}, \emph{pratinidhisu}.) Politics? (\textbf{\kn{ರಾಜಕೀಯ}}, \emph{rājakīya}.) A politician? (\textbf{\kn{ರಾಜಕಾರಣಿ}}, \emph{rājakāraṇi}.) A political party? (\textbf{\kn{ಪಕ್ಷ}}, \emph{pakṣa}.) To boycott? (\textbf{\kn{ಬಹಿಷ್ಕರಿಸು}}, \emph{bahiṣkarisu}.) Leadership? (\textbf{\kn{ನಾಯಕತ್ವ}}, \emph{nāyakatva}.) Administration? (\textbf{\kn{ಆಡಳಿತ}}, \emph{āḍaḷita}.) A ministry? (\textbf{\kn{ಸಚಿವಾಲಯ}}, \emph{sacivālaya}.) A village council? (\textbf{\kn{ಪಂಚಾಯಿತಿ}}, \emph{pañcāyiti}.) To implement? (\textbf{\kn{ಜಾರಿಗೊಳಿಸು}}, \emph{jārigoḷisu}.) A committee? (\textbf{\kn{ಸಮಿತಿ}}, \emph{samiti}.) A draft law? (\textbf{\kn{ಮಸೂದೆ}}, \emph{masūde}.) To approve? (\textbf{\kn{ಅಂಗೀಕರಿಸು}}, \emph{aṅgīkarisu}.) An amendment? (\textbf{\kn{ತಿದ್ದುಪಡಿ}}, \emph{tiddupaḍi}.) A legislator? (\textbf{\kn{ಶಾಸಕ}}, \emph{śāsaka}.) To endorse? (\textbf{\kn{ಅನುಮೋದಿಸು}}, \emph{anumōdisu}.) A citizen? (\textbf{\kn{ನಾಗರಿಕ}}, \emph{nāgarika}.) Citizenship? (\textbf{\kn{ಪೌರತ್ವ}}, \emph{pauratva}.) Public? (\textbf{\kn{ಸಾರ್ವಜನಿಕ}}, \emph{sārvajanika}.) A community? (\textbf{\kn{ಸಮುದಾಯ}}, \emph{samudāya}.) To unite? (\textbf{\kn{ಒಗ್ಗೂಡು}}, \emph{oggūḍu}.)
\end{tcolorbox}
